\section*{Capítulo \ref{ch:b2:mantle-melting} --- A fusão do manto}

\begin{solution}{exo:b2:mantle-melting:1}
(a) Mais de \qty{90}{\percent} de olivina: \omterm{def:b2:mantle-melting:peridotites}{dunito}. (b) Olivina e ortopiroxênio
sem clinopiroxênio: \omterm{def:b2:mantle-melting:peridotites}{harzburgito}. (c) Mais de \qty{40}{\percent} de olivina com
os dois piroxênios: \omterm{def:b2:mantle-melting:peridotites}{lherzolito}. (d) Menos de \qty{40}{\percent} de olivina: um piroxenito
(um olivina-websterito, com os dois piroxênios). (e) $57/98 = \qty{58.2}{\percent}$,
$28/98 = \qty{28.6}{\percent}$, $13/98 = \qty{13.3}{\percent}$: um \omterm{def:b2:mantle-melting:peridotites}{lherzolito}.
\end{solution}

\begin{solution}{exo:b2:mantle-melting:2}
(a) O liquidus a \qty{60}{\percent} de anortita é de \qty{1390}{\celsius}, no
campo da anortita: a anortita cristaliza primeiro. (b) O diopsídio aparece no
eutético, a \qty{1270}{\celsius}, quando o líquido tem \qty{42}{\percent} de anortita.
(c) O primeiro líquido de qualquer mistura dos dois minerais é o líquido eutético,
com \qty{42}{\percent} de anortita, a \qty{1270}{\celsius}.
\end{solution}

\begin{solution}{exo:b2:mantle-melting:3}
(a) $\num{4.0e-5}\times1623/(3300\times1000) = \qty{1.97e-8}{K/Pa} =
\qty{19.7}{K/GPa}$. (b) $\num{e9}/(3300\times9.81) = \qty{3.09e4}{m}$, ou
\qty{30.9}{km} por gigapascal; $19.7/30.9 = \qty{0.64}{K/km}$, cerca de cinquenta vezes
menos que os \qty{30}{K/km} perto da superfície, onde o calor é conduzido, e não
transportado.
\end{solution}

\begin{solution}{exo:b2:mantle-melting:4}
(a) $T_{\mathrm s}(0) = \qty{1085.7}{\celsius}$; $T_{\mathrm s}(1) = 1085.7 +
132.9 - 5.1 = \qty{1213.5}{\celsius}$; $T_{\mathrm s}(2) = 1085.7 + 265.8 - 20.4 =
\qty{1331.1}{\celsius}$; $T_{\mathrm s}(3) = 1085.7 + 398.7 - 45.9 =
\qty{1438.5}{\celsius}$. (b) $132.9 - 2\times5.1\times2 = \qty{112.5}{K/GPa}$.
(c) A \qty{3}{GPa}, \qty{1400}{\celsius} está abaixo do solidus,
\qty{1438.5}{\celsius}; subindo até \qty{2}{GPa} ao longo da adiabática, o manto esfria apenas
cerca de \qty{20}{\celsius}, até uns \qty{1380}{\celsius}, enquanto o solidus cai para
\qty{1331.1}{\celsius}: a rocha está agora acima de seu solidus e parte dela funde.
\end{solution}

\begin{solution}{exo:b2:mantle-melting:5}
(a) $1350 + (1277.5 - 1350)\times(30 - 20)/(40 - 20) = \qty{1314}{\celsius}$.
(b) A composição fica do lado do diopsídio em relação ao eutético: diopsídio.
(c) O líquido atinge o eutético com \qty{42}{\percent} de anortita; a anortita
permanece no líquido, de modo que a fração de líquido é $30/42 = 0.714$ e $1 - 0.714 =
\qty{29}{\percent}$ do sistema cristalizou como diopsídio. (d) Tudo
cristaliza: \qty{70}{\percent} de diopsídio e \qty{30}{\percent} de anortita, a
composição global.
\end{solution}

\begin{solution}{exo:b2:mantle-melting:6}
(a) Forsterita: $60.08/(2\times40.30 + 60.08) = 60.08/140.68 =
\qty{42.71}{\percent}$ de \ce{SiO2}; enstatita: $60.08/100.38 = \qty{59.85}{\percent}$.
(b) $0.975\times59.85 + 2.5 = \qty{60.85}{\percent}$. (c) A enstatita fica entre a
forsterita e o líquido: a fração de forsterita é $(60.85 - 59.85)/(60.85 -
42.71) = 1.00/18.14 = \qty{5.5}{\percent}$, o valor de Bowen e Andersen. (d) O
fundido é mais rico em sílica que o ortopiroxênio de onde vem, e o resíduo
ganha olivina.
\end{solution}

\begin{solution}{exo:b2:mantle-melting:7}
(a) $5.1P^2 - 112.9P + 214.3 = 0$: $P_0 = (112.9 - \sqrt{12746.4 -
4371.7})/10.2 = (112.9 - 91.51)/10.2 = \qty{2.10}{GPa}$, a $2.097\times30.9 =
\qty{64.8}{km}$. (b) $5.1P^2 - 112.9P + 364.3 = 0$: $P_0 = (112.9 -
72.90)/10.2 = \qty{3.92}{GPa}$, a \qty{121}{km}. (c) $T_{\mathrm s}'(2.10) = 132.9
- 21.4 = \qty{111.5}{K/GPa}$, logo $\Delta P_0 = 10/(111.5 - 20) = \qty{0.109}{GPa}$
por \qty{10}{\celsius}, ou \qty{3.4}{km}.
\end{solution}

\begin{solution}{exo:b2:mantle-melting:8}
(a) $T_{\mathrm s} = \qty{1331.1}{\celsius}$, $T_{\mathrm l} = 1475.0 + 160.0 -
12.8 = \qty{1622.2}{\celsius}$; $T = 1331.1 + 0.10\times291.1 =
\qty{1360.2}{\celsius}$, ou \qty{1633}{K}. (b) $(\partial T/\partial P)_F = 112.5
+ 0.10\times(67.2 - 112.5) = \qty{108.0}{K/GPa}$; $(\partial T/\partial F)_P =
T_{\mathrm l} - T_{\mathrm s} = \qty{291.1}{K}$. (c) $c_p/T = 0.612$; com $\alpha/\rho
= \qty{12.1}{J/(kg.K)}$ por gigapascal, o numerador é $0.612\times108.0 - 12.1
= 54.0$ e o denominador $300 + 0.612\times291.1 = 478.2$: $-\mathrm dF/\mathrm
dP = 0.113$, \qty{11.3}{\percent} por gigapascal. (d) $54.0/178.2 = 0.303$,
\qty{30}{\percent} por gigapascal: o calor latente reduz a produtividade por um
fator 2.7.
\end{solution}

\begin{solution}{exo:b2:mantle-melting:9}
(a) $43\times1 = \qty{43}{\celsius}$; $43\times5^{0.75} = 43\times3.34 =
\qty{144}{\celsius}$; $43\times13^{0.75} = 43\times6.85 = \qty{294}{\celsius}$.
(b) $1213.5 - 294 = \qty{919}{\celsius}$. (c) \qty{125}{ppm} são
\qty{0.0125}{\percent}, logo $X = 0.0125/0.01 = \qty{1.25}{\percent}$ e $\Delta T =
43\times1.25^{0.75} = 43\times1.18 = \qty{51}{\celsius}$.
\end{solution}

\begin{solution}{exo:b2:mantle-melting:10}
(a) A base tem largura $2H$ e o manto a atravessa a $U$: $2UH$ por unidade de comprimento
de dorsal e de tempo. Cada lado abrange a altura $H$ e é atravessado horizontalmente a
$U$: $UH$ cada um, $2UH$ no total. Com lados a $\theta$, a base teria largura
$2H/\tan\theta$ e o fluxo de entrada seria $2UH/\tan\theta$, igual ao de saída somente se
$\tan\theta = 1$: nesse modelo, o ângulo é fixado pela conservação da massa. (b) Com
produtividade constante, a \omterm{def:b2:mantle-melting:fraction}{fração de fundido} no vértice é $F_{\max} = \Pi P_0 = \Pi\rho
g z_0$, e $H = z_0$: $h_c = F_{\max}z_0/2 = \Pi\rho g z_0^2/2$. (c) A cada ano,
\qty{50}{mm} de placa nova, $\qty{0.05}{m}\times\qty{1000}{m}\times\qty{7100}{m} =
\qty{3.6e5}{m^3}$ de crosta por quilômetro de dorsal; com \qty{12}{\percent} de \omterm{def:b2:mantle-melting:fraction}{fração
de fundido} média, ela vem de $\num{3.6e5}/0.12 = \qty{3.0e6}{m^3}$ de manto.
\end{solution}

\begin{solution}{exo:b2:mantle-melting:11}
(a) Depois que uma massa $t$ de A deixou uma unidade de massa de líquido, o líquido contém $0.6 -
t$ de A, 0.2 de B e 0.2 de C em uma massa $1 - t$: $\big((0.6 - t), 0.2,
0.2\big)/(1 - t)$, e B/C permanece 1. (b) Sobre a cotética, $(0.40 - 0.15s, 0.60 -
0.30s, 0.45s)$; B = C dá $0.60 - 0.30s = 0.45s$, $s = 0.8$: $L_1$ tem A 28, B
36, C 36. Então $0.2/(1 - t) = 0.36$, $t = 0.44$: \qty{44}{\percent} do sistema
cristalizou como A. (c) Todo o C está no líquido: $f = 20/45 = 0.44$,
\qty{44}{\percent} ainda líquidos. (d) A: $0.60 - 0.444\times0.25 = 0.49$; B: $0.20 -
0.444\times0.30 = 0.07$; com o líquido, $49 + 7 + 44 = \qty{100}{\percent}$.
\end{solution}

\begin{solution}{exo:b2:mantle-melting:12}
(a) $X_{\mathrm{sat}}(1) = 12 + 1 = \qty{13}{\percent}$, $\Delta T =
\qty{294}{\celsius}$, solidus saturado $1213.5 - 294 = \qty{919}{\celsius}$;
$X_{\mathrm{sat}}(2) = 12\times1.516 + 2 = \qty{20.2}{\percent}$, $\Delta T =
43\times20.2^{0.75} = 43\times9.53 = \qty{410}{\celsius}$, solidus saturado $1331.1
- 410 = \qty{921}{\celsius}$. (b) O fundido dissolve mais água a pressão mais alta,
e o rebaixamento maior compensa a subida de \qty{118}{\celsius} do solidus seco.
(c) O solidus precisa cair $1331.1 - 1100 = \qty{231}{\celsius}$: $X =
(231/43)^{1/0.75} = 5.37^{1.333} = \qty{9.4}{\percent}$ de água, menos que os
\qty{20.2}{\percent} da saturação, de modo que a cunha pode fundir.
\end{solution}

\begin{solution}{pb:b2:mantle-melting:1}
\textbf{1.} $1350 + 273 = \qty{1623}{K}$.
\textbf{2.} $\num{4.0e-5}\times1623/(3300\times1000) = \qty{1.97e-8}{K/Pa} =
\qty{19.7}{K/GPa}$, arredondado para \qty{20}{K/GPa}.
\textbf{3.} $\num{e9}/(3300\times9.81) = \qty{30.9}{km}$; $20/30.9 =
\qty{0.65}{K/km}$.
\textbf{4.} $1350 + 20P = 1085.7 + 132.9P - 5.1P^2$, isto é, $5.1P^2 - 112.9P +
264.3 = 0$: $P_0 = (112.9 - \sqrt{12746.4 - 5391.7})/10.2 = (112.9 - 85.76)/10.2 =
\qty{2.66}{GPa}$.
\textbf{5.} $2.66\times30.9 = \qty{82}{km}$; $T = 1350 + 20\times2.66 =
\qty{1403}{\celsius}$.
\textbf{6.} $1085.7 + 132.9\times2.66 - 5.1\times2.66^2 = 1085.7 + 353.5 - 36.1 =
\qty{1403}{\celsius}$.
\textbf{7.} $T_{\mathrm s} = 1085.7 + 199.4 - 11.5 = \qty{1273.6}{\celsius}$;
$T_{\mathrm l} = 1475.0 + 120.0 - 7.2 = \qty{1587.8}{\celsius}$; diferença de
\qty{314.2}{K}.
\textbf{8.} $1273.6 + 0.10\times314.2 = \qty{1305.0}{\celsius}$, \qty{1578}{K}.
\textbf{9.} $T_{\mathrm s}' = 132.9 - 15.3 = \qty{117.6}{K/GPa}$, $T_{\mathrm l}' =
80.0 - 9.6 = \qty{70.4}{K/GPa}$: $(\partial T/\partial P)_F = 117.6 +
0.10\times(70.4 - 117.6) = \qty{112.9}{K/GPa}$; $(\partial T/\partial F)_P =
\qty{314.2}{K}$.
\textbf{10.} $c_p/T = 1000/1578 = 0.634$; numerador $0.634\times112.9 - 12.1 = 71.6
- 12.1 = 59.5$; denominador $300 + 0.634\times314.2 = 499.2$: $0.119$, ou
\qty{11.9}{\percent} por gigapascal.
\textbf{11.} Com a função ajustada $F = T'^{1.5}$, sendo $T'$ a fração do caminho entre o solidus e o liquidus, a temperatura precisa subir abruptamente com $F$ para produzir
o primeiro por cento de fundido, de modo que $(\partial T/\partial F)_P$ é grande perto do solidus
e a fusão começa devagar. $F_{\max} = 0.09\times2.66 = 0.24$, \qty{24}{\percent}.
\textbf{12.} $0.13/0.50 = 0.26$: o clinopiroxênio se esgota a \qty{26}{\percent} de
fusão. A \qty{24}{\percent}, o resíduo sob o eixo ainda contém um pouco dele:
um \omterm{def:b2:mantle-melting:peridotites}{lherzolito} pobre em clinopiroxênio, no limiar do \omterm{def:b2:mantle-melting:peridotites}{harzburgito}.
\textbf{13.} $h_c = 0.24\times82/2 = \qty{9.8}{km}$.
\textbf{14.} $F_{\max}/2 = \qty{12}{\percent}$.
\textbf{15.} \qty{9.8}{km} são \qty{38}{\percent} a mais que \qty{7.1}{km}, acima do
intervalo da crosta normal.
\textbf{16.} A produtividade constante dá fundido demais ao manto mais profundo; o
modelo ajustado produz o primeiro fundido devagar, $F(z)$ permanece mais baixo perto do início, e
$\int F\,\mathrm dz$ é menor.
\textbf{17.} $F_{\max} = \Pi P_0$ e $H = 30.9P_0$ (km): $h_c = \Pi P_0\times30.9P_0/2
= kP_0^2$ com $k = 0.09\times30.9/2 = \qty{1.39}{km/GPa^2}$.
\textbf{18.} $P_0 = \sqrt{7.1/1.39} = \sqrt{5.11} = \qty{2.26}{GPa}$.
\textbf{19.} $2.26\times30.9 = \qty{70}{km}$.
\textbf{20.} $T_{\mathrm s}(2.26) = 1085.7 + 300.4 - 26.0 = \qty{1360.0}{\celsius}$,
e $T_p = 1360.0 - 20\times2.26 = 1360.0 - 45.2 = \qty{1315}{\celsius}$.
\textbf{21.} Para \qty{6.3}{km}: $P_0 = \sqrt{4.53} = \qty{2.13}{GPa}$, $T_{\mathrm s}
= \qty{1345.6}{\celsius}$, $T_p = \qty{1303}{\celsius}$. Para \qty{7.9}{km}: $P_0 =
\sqrt{5.68} = \qty{2.38}{GPa}$, $T_{\mathrm s} = \qty{1373.1}{\celsius}$, $T_p =
\qty{1326}{\celsius}$. A dispersão da crosta permite $1315 \pm
\qty{12}{\celsius}$.
\textbf{22.} $P_0 = \sqrt{20/1.39} = \sqrt{14.39} = \qty{3.79}{GPa}$; $T_{\mathrm s} =
1085.7 + 503.7 - 73.3 = \qty{1516.1}{\celsius}$; $T_p = 1516.1 - 75.8 =
\qty{1440}{\celsius}$, um excesso de $1440 - 1315 = \qty{125}{\celsius}$.
\textbf{23.} As duas usam evidências diferentes. As lavas registram o manto mais quente,
no centro da pluma; a crosta faz a média dos fundidos de todo o manto que
atravessa a região de fusão, boa parte dele mais frio. O modelo simples também mantém
a produtividade constante, embora o clinopiroxênio se esgote a \qty{26}{\percent},
ao passo que aqui $\Pi P_0 = 0.09\times3.79 = 0.34$: o manto quente real funde menos, e precisa
de uma temperatura mais alta para produzir a mesma crosta.
\textbf{24.} $T_{\mathrm s}'(2.26) = 132.9 - 23.1 = \qty{109.8}{K/GPa}$; $\Delta P_0 =
10/(109.8 - 20) = \qty{0.111}{GPa}$; como $h_c \propto P_0^2$, $\Delta h_c/h_c =
2\times0.111/2.26 = \qty{9.8}{\percent}$, cerca de um décimo por \qty{10}{\celsius}.
\textbf{25.} Em uma dorsal muito lenta, o manto sobe tão devagar que perde calor por
condução em direção à superfície fria; ele deixa a adiabática, sua fusão para
mais fundo, e o topo do triângulo, onde $F$ é máximo, é cortado: a crosta é
mais delgada.
\textbf{26.} A crosta das dorsais normais dá uma \omterm{def:b2:mantle-melting:tp}{temperatura potencial do manto} de
\textbf{\qty[text-series-to-math=true]{1315}{\celsius}}, a cerca de
\qty{12}{\celsius}: \qty{35}{\celsius} abaixo dos \qty{1350}{\celsius} da física
mineral, que é incerta em aproximadamente o mesmo tanto, e bem dentro do intervalo de 1280 a
\qty{1400}{\celsius} das lavas. Três termômetros independentes concordam em um manto
perto de \qty{1300}{\celsius} a \qty{1350}{\celsius}.
\end{solution}
